% === Begin Label Data ===
% ID: 209
% 难度星级: 
% 标签: 
% 备注: 
% 组卷引用次数: 0
% === End  Label Data ===

\begin{problem}{2026}{M}{武汉市2026届高中毕业生三月调研考试}{19}{数列}
有 $n$ 张编号分别为 $1$ 到 $n$ 的卡片，横向随机排列。对于这 $n$ 张卡片，初始状态下单片标号从左到右为 $A_1, A_2, \cdots, A_n$，记此时的卡片排列为 $(A_1, A_2, \cdots, A_n)$。对这 $n$ 张卡片的排列进行如下三步操作：  
1. 取出最左边的卡片，记其标号为 $k$；  
2. 剩余卡片中，标号小于 $k$ 的卡片按照原排列中的从左到右顺序依次为 $L_1, L_2, \cdots, L_{k-1}$（若不存在则为空），标号大于 $k$ 的卡片按照原排列中的从左到右顺序依次为 $R_1, R_2, \cdots, R_{n-k}$（若不存在则为空）；  
3. 对这 $n$ 张卡片重新排列，得到新排列：$(L_1, L_2, \cdots, L_{k-1}, k, R_1, R_2, \cdots, R_{n-k})$。  
每进行完上述三步操作，称为一次“完整操作”。

（1）若初始排列为 $(3,5,2,4,1)$，写出连续经过两次完整操作后得到的新排列；

（2）求初始排列经过一次完整操作后恰好能得到 $(1,2,\cdots,n)$ 的顺序排列的概率；

（3）记初始排列中有 $B_n$ 个排列种数能经过连续若干次完整操作后能得到 $(1,2,\cdots,n)$ 的顺序排列，当 $n\geq 2$ 时，证明：$B_{n+1}\leq nB_n + B_{n-1}$。
\end{problem}